\documentclass[11pt]{article}
\usepackage{fullpage,graphicx,hyperref,amssymb,amsmath}
\newcommand{\fig}[2]{\begin{figure}[h]\begin{center}%
  \includegraphics[scale=#2]{#1.pdf}\end{center}%
  \end{figure}}
\newcommand{\down}[1]{\left\lfloor #1\right\rfloor}
\newcommand{\Ex}{\mathbb{E}}

\begin{document}

\begin{center}\Large\bf 
CS 473 (Spring 2025)\\
Homework 4 (due Mar 6 Thu 10am)
\end{center}

\medskip\noindent
{\bf Instructions:} As previous homeworks.

\begin{description}
\bigskip
\item[Problem 4.1:]
Suppose we have a Las Vegas randomized algorithm A that 
runs in $\mu$ expected number of steps.  As mentioned in class, for any given $t>1$,
we can obtain a new Monte Carlo algorithm A$'$ that runs within $t\mu$
steps by outputting ``garbage'' if the time limit is exceeded.
The error probability is at most $1/t$ by Markov's inequality.

But we can do better!  For example, we can run A for up to $2\mu$ steps,
and if it fails, rerun A for another $2\mu$ steps (with an independent
sequence of random bits), etc., for $\down{t/2}$ rounds.
This gives a Monte Carlo algorithm with
the total number of steps still bounded by $t\mu$, and the error probability reduced to $1/2^{\down{t/2}} = O(1/2^{t/2}) \le O(0.708^t)$,
which decays exponentially!

In this question, you will analyze further improvements to the base of the exponential:

\begin{enumerate}
\item[(a)] (25 pts) 
Let $X\ge 0$ be a random variable with $\Ex[X]=\mu$.
Given $\alpha_1,\alpha_2,\ldots,\alpha_k\ge 0$, let 
$p_1=\Pr[X\ge \alpha_1\mu]$, $p_2=\Pr[X\ge (\alpha_1+\alpha_2)\mu]$, \ldots,
$p_k=\Pr[X\ge (\alpha_1+\cdots+\alpha_k)\mu]$.

Prove the following inequality:
\[ \alpha_1 p_1 + \alpha_2 p_2 + \cdots + \alpha_k p_k\ \le\ 1. 
\]

[Hint: modify the proof of Markov's inequality and use telescoping sum\ldots]

\smallskip
\item[(b)] (25 pts) Now, consider a strategy with three iterations 
where we run A for $\alpha_1\mu$ steps in the first iteration, $(\alpha_1+\alpha_2)\mu$ steps in the second iteration,
and $(\alpha_1+\alpha_2+\alpha_3)\mu$ steps in the third iteration.

With this strategy, we obtain a Monte Carlo algorithm with
the total number of steps upper-bounded
by $(3\alpha_1 + 2\alpha_2 + \alpha_3)\mu$.  Bound the error probability
purely as a function of $\alpha_1$, $\alpha_2$, and $\alpha_3$.

[Hint: you may use the fact that $x+y+z\le 1\ (x,y,z\ge 0)$ implies $xyz\le 1/27$.  (This
follows from the standard inequality that the ``geometric mean'' is at most the ``arithmetic mean''.)]

\smallskip
\item[(c)] (25 pts) Using part (b), show how to obtain a Monte Carlo algorithm with
the total number of steps bounded by $t\mu$, and the error probability 
at most $O(0.545^t)$ (thus, improving $O(0.708^t)$).

[Hint: repeat for 
$\down{\frac{t}{3\alpha_1 + 2\alpha_2 + \alpha_3}}$ rounds; set $\alpha_1=c/3$, $\alpha_2=c/2$, $\alpha_3=c$ for some choice of $c$ (you might need a calculator)\ldots]

\smallskip
\item[(d)] (25 pts) Generalize part (b) from 3 iterations to $k$ iterations for a larger constant $k$, and show how to reduce the error probability further to below $O(0.4^t)$.

[Hint: choose some $k$ below 50, and choose $c=1$\ldots]
\end{enumerate}


\newpage
\textbf{Solution:}
\begin{enumerate}
  \item[(a)] Assume $\beta_i = (\alpha_1+\ldots+\alpha_i)\mu, \beta_0 = 0$,
  \begin{align*}
  \mathbb{E}[X] 
  &= \int_{0}^{\infty}\Pr[X\ge t] dt \\
  &= \sum_{i=1}^{k} \int_{\beta_{i-1}}^{\beta_{i}} \Pr[X\ge t]dt 
  + \int_{\beta_k}^{\infty} \Pr[X\ge t]dt
  \end{align*}
  Since $\Pr[X\ge t]$ is non-increase in $t$,
  \begin{align*}
  \int_{\beta_{i-1}}^{\beta_i} \Pr[X\ge t]dt 
  &\le Pr[X\ge \beta_i](\beta_i - \beta_{i-1}) \\
  &= p_i\alpha_i\mu
  \end{align*}
  Thus,
  \[
  \mathbb{E}[X] \le \sum_{i=1}^{k} p_i\alpha_i\mu
  \]\[
  \sum_{i=1}^{k} \alpha_i p_i \le 1
  \]

  \item[(b)] The probability that the Monte Carlo algorithm fails is 
  $p_1 p_2 p_3$, from (a),
  \[
  \alpha_1p_1 + \alpha_2p_2 + \alpha_3p_3 \le 1 \ 
  (\alpha_1p_1, \alpha_2p_2, \alpha_3p_3 \ge 0)
  \]
  Thus,
  \[
  \alpha_1p_1\alpha_2p_2\alpha_3p_3 \le \frac{1}{27}
  \]\[
  \Pr(\text{error}) = p_1 p_2 p_3 \le \frac{1}{27\alpha_1\alpha_2\alpha_3}
  \]

  \item[(c)] Let $\alpha_1=\frac{c}{3},\alpha_2=\frac{c}{2},\alpha_3=\frac{c}{1}$,
  the Monte Carlo algorithm repeats for $\down{\frac{t}{3c}}$ rounds, each
  round takes $3c\mu$ steps. The total number of steps is bounded by,
  \[\down{\frac{t}{3c}} \cdot 3c\mu \le t\mu\]
  The probability that the algorithm fails is,
  \[\Pr(\text{error}) \le (\frac{6}{27c^3})^{\down{\frac{t}{3c}}} 
  = O((\frac{6}{27c^3})^{\frac{t}{3c}})\]
  Let $c=\frac{\sqrt[3]{6}e}{3}\approx1.646$, we get,
  \[\Pr(\text{error}) \le O(0.545^t)\]

  \item[(d)] Let $\alpha_i=\frac{1}{k+1-i}$ for $i=1,\ldots,k$.
  The Monte Carlo algorithm repeats for $\down{\frac{t}{k}}$ rounds, each
  round takes $k\mu$ steps. The total number of steps is bounded by $t\mu$.
  The probability that each round fails is,
  \[
  \prod_{i=1}^{k}p_i \le k^{-k}\prod_{i=1}^{k}\alpha_i^{-1} 
  = k^{-k}k!
  \]
  The probability that the algorithm fails is,
  \[
  \Pr(\text{error}) \le (k^{-k}k!)^{\down{\frac{t}{k}}} = O((k^{-k}k!)^{\frac{t}{k}})
  \]
  $\Pr(\text{error})$ is decresing when $k>0$, let $k\ge32$,
  \[\Pr(\text{error}) \le O(0.3997^t) \le O(0.4^t)\]
\end{enumerate}



\newpage
\item[Problem 4.2:]
In a popular form of logic puzzles, you 
land on an island that has three types of inhabitants:
``knights'', who always tell the truth;
``knaves'', who always lie; and ``spies'', who sometimes
lie and sometimes tell the truth.  

Suppose there are $n$ inhabitants, where 60\% are known to be decent
folks, i.e., knights.  The remaining 40\% are bad, i.e., knaves or
 spies.  
You want to know who are the good/bad guys, i.e., 
determine the types of all $n$ inhabitants.  You are allowed to ask 
only questions of the form, ``is person A a knight/knave/spy?'', 
to another person B\@.  
(All the inhabitants know each other.)
Obviously, if you can find a person who you
know is a knight, the problem can be solved after asking 
$n$ additional questions.
\begin{enumerate}
\item[(a)] (10 pts) Give a (very) efficient Monte Carlo algorithm that finds a knight.  
State the probability of error.  (Hint: this is supposed to be very easy!)

\smallskip
\item[(b)] (90 pts) Give a Las Vegas algorithm that finds a knight by asking
$O(n)$ expected number of questions.  Analyze the constant factor 
in
the big-Oh and make it smaller than 1.5.

[Hint: use (a).  How can you confirm whether a specific person
is a knight by asking $O(n)$ questions?]

[Note: there is also a deterministic algorithm that requires 
$O(n)$ questions, but it is more complicated and has a larger constant factor.]
\end{enumerate}

\textbf{Solution:}
\begin{enumerate}
  \item[(a)] Monte Carlo algorithm: \\
  Randomly select a person $A$ and accept $A$ as a knight.

  The probability of error is 0.4.

  \item[(b)] Las Vegas algorithm: \\
  1. Select a random person $A$. \\
  2. Select a group of $0.8n$ inhabitants and ask each one of them whether
  $A$ is a knight. \\
  3. $A$ is a knight if more than half of the group agree since there are
  at most 0.4n knaves. \\
  4. Otherwise, repeat step $1\sim3$ until we find a knight.

  The expected number of questions is,
  \begin{align*}
  \mathbb{E}(questions) 
  &= \sum_{i=1}^{\infty} i\cdot(0.4)^{i-1}\cdot0.6\cdot0.8n \\
  &= \sum_{i=0}^{\infty} (i+1)\cdot(0.4)^i\cdot0.48n \\
  &= \left(\frac{0.4}{(1-0.4)^2} + \frac{1}{1-0.4}\right) \cdot 0.48n \\
  &\approx 1.333n < 1.5n
  \end{align*}
\end{enumerate}


\newpage
\item[Problem 4.3:]
We are given a set $S$ of strings of total length $N$ over the binary alphabet $\{0,1\}$.
We are also given a number $\ell$. 
We want to find a palindrome $p$ of length $\ell$ and two distinct strings $s,s'\in S$ such that
their concatenation $ss'$ contains $p$ as a substring.

For example, for $S=\{0111,1111010,10110,11011\}$ and $\ell=7$, 
the concatenation of $1111010$ and $11011$ contains the palindrome $1101011$.

Describe a randomized Las Vegas algorithm for this problem, by using the Karp--Rabin fingerprint technique (to be covered in Feb 28's lecture).  The expected running time should be $O(N\log N)$ or better.

[Hint: for every string $s\in S$, pre-compute fingerprints of each prefix/suffix of $s$ and the reverse of $s$.  You may find subtraction useful\ldots]

\textbf{Solution:}

1. For every string $s\in S$, pre-compute Karp-Rabin fingerprints of each prefix/suffix 
$x$ of $s$ and the reverse of $s$ as $F(x)$ and $F(x^R)$.

2. For every string $s\in S$, use a slide window of length $\ell$ to find whether
there is a palindrome within $s$.

For each suffix $x$ of length $i$ and prefix $y$ of length $\ell-i$, $xy$ might be a palindrome if
\[F(x)\cdot b^{\ell-i} + F(y) = F(y^R)\cdot b^i + F(x^R)\]
where $b$ is the base in Karp-Rabin fingerprints.

3. Sort all prefix $y$ according to $F(y^R)\cdot b^i - F(y)$. 

For each suffix $x$,
compute $F(x)\cdot b^{\ell-i} - F(x^R)$ and check if there is a corresponding $y$ using binary search.

4. If there is a possible match, verify if it is a palindrome. 

If not, find palindrome by brute force.

5. Repeat step 3 and 4 until the palindrome is found.

Runtime Analysis:

$S$ has at most $\frac{2N}{\ell}$ strings of length over $\frac{\ell}{2}$
(otherwise, there won't be any palindrome of length $\ell$), 
each string has at most $\ell-1$ prefix/suffix, total number of prefix/suffix is $O(N)$, 
by using subtraction, all Karp-Rabin fingerprints can be computed in $O(N)$ time. 
Step 1 and 2 take $O(N)$ time, step 3 takes $O(N\log N)$time.
Let $p=N^{d+2}$ in fingerprints, the probability of using brute force in step 5 is
$O(\frac{1}{N^d})$, expected time complexity is $O(\ell N \times \frac{1}{N^d})$,
which is smaller than $O(N)$ by setting $d=1$.
Total time complexity is expected to be $O(N\log N)$.

\end{description}
\end{document}
